\chapter{Struttura atomica e modelli nucleari}
\section{Modello atomico e sezione d'urto}
Il primo modello atomico fu proposto da Dalton nel 1803, che immaginò gli atomi come sfere indivisibili e indistruttibili. Successivamente, nel 1897, Thomson scoprì l'elettrone e propose il modello a "panettone", in cui gli elettroni erano immersi in una massa positiva. Nel 1911, Rutherford condusse l'esperimento di scattering delle particelle alfa, che portò alla scoperta del nucleo atomico e alla proposta del modello planetario, dove gli elettroni orbitano attorno a un nucleo centrale carico positivamente.

\addtikz[Modello di Rutherford]{
    \draw[thick] (0,0) circle (2);
    \draw[thick] (0,0) circle (3);
    \fill[thick,red!50,opacity=0.5] (0,0) circle (0.4);
    \draw[thick,red] (0,0) circle (0.4)node {+};
    \fill[supportDarkBlue] (1.42,1.42) circle (0.1)node[above right] {$e^-$};
    \fill[supportDarkBlue] (-1.42,-1.42) circle (0.1)node[above right] {$e^-$};
    \fill[supportDarkBlue] (3,0) circle (0.1)node[above right] {$e^-$};
    \fill[supportDarkBlue] (0,3) circle (0.1)node[above right] {$e^-$};
    \fill[supportDarkBlue] (-3,0) circle (0.1)node[above right] {$e^-$};
    \fill[supportDarkBlue] (0,-3) circle (0.1)node[above right] {$e^-$};
}{1}{Rutherford_model}

Durante l'esperimento si possono registrare angoli di riflessione molto ampi, fino a 180°. Questo fatto è incompatibile con il modello di Thomson, che prevede deviazioni molto piccole.\\
Rutherford propose quindi un modello in cui la maggior parte della massa e della carica positiva dell'atomo è concentrata in un nucleo centrale molto piccolo, mentre gli elettroni orbitano attorno ad esso a distanze relativamente grandi.
Infine, nel 1913, Bohr introdusse il concetto di quantizzazione delle orbite elettroniche, spiegando la stabilità degli atomi e le righe spettrali osservate.\\
Utilizzando il modello geometrico (di Rutherford) ipotizziamo di avere un flusso di particelle monocromatico, ovvero con una sola energia, che colpisce perpendicolarmente un bersaglio sottile. Abbiamo:
\begin{itemize}
 \item $N_b$ centri di diffusione (o scattering)
 \item $\Phi_a$ flusso di particelle
 \item $\sigma_b$ sezione d'urto del bersaglio
\end{itemize}
Otteniamo dunque un tasso di interazioni per unità di volume pari a:
\begin{equation}
 \dv{N}{t} = \Phi_a N_b \sigma_b
\end{equation}
Possiamo dunque dare una definizione di sezione d'urto:
\begin{definition}{Sezione d'urto}
    La \textbf{sezione d'urto} $\sigma$ è definita come il rapporto tra il numero di interazioni per unità di tempo e il flusso di particelle incidente moltiplicato per il numero di centri di diffusione:
    \begin{equation}
     \sigma = \dfrac{1}{\Phi_a N_b} \dv{N}{t}
    \end{equation}
\end{definition}
La sezione d'urto è la grandezza che caratterizza le singole interazioni. Ad esempio la fissione su $^{235}\text{U}$ ha una sezione d'urto variabile nell'ordine dei barn, dove $\qty{1}{\barn} = \qty{e-28}{\meter} = \qty{100}{\femto\meter\squared}$.

\addfigure[Sezione d'urto per la fissione del $^{235}\text{U}$ in funzione dell'energia del neutrone incidente. Fonte: \href{https://whatisnuclear.com/barn-jams.html}{What is nuclear?}]{png/Cross_section_U235}{0.7}

All'aumento della sezione d'urto corrisponde un aumento del numero di interazioni, dunque per reazioni di fissione dobbiamo lavorare a basse energie.\\
Possiamo anche definire una sezione d'urto differenziale.
\begin{definition}{Sezione d'urto differenziale}
    La \textbf{sezione d'urto differenziale} è definita come la variazione infinitesima della sezione d'urto in funzione dell'angolo solido:
    \begin{equation}
     \dv{\sigma}{\Omega} = \lim_{\Delta \Omega \to 0} \dfrac{\Delta \sigma}{\Delta \Omega}
    \end{equation}
\end{definition}
Con $\sigma$ definito in coordinate polari:
\begin{equation}
 \sigma = \int_{4 \pi} \dv{\sigma}{\Omega} \dd{\Omega} = \int_0^{2\pi} \dd{\varphi} \int_{-1}^{1} \dd{\brackets{\cos \theta}} \dv{\sigma}{\Omega}
\end{equation}
Per piccole sezioni d'urto vale la proporzionalità:
\begin{equation}
 \sigma = \dv{\sigma}{\Omega} \Delta \Omega
\end{equation}
Dove $\Delta \Omega$ è l'angolo solido sotteso dalla sezione d'urto.\\
Poiché vogliamo attraversare una regione di spazio con un potenziale, dobbiamo considerare elettroni con un'energia dell'ordine di qualche \si{\mega\electronvolt}. Infatti, per una barriera di circa $\qty{1}{\femto\meter}$, utilizzando la relazione di De Broglie:
\begin{equation}
 \dfrac{\lambda}{2 \pi} = \dfrac{\hbar c}{p c} = \dfrac{\hbar}{K_ec} \simeq \dfrac{\qty{200}{\mega\electronvolt\femto\meter}}{K_e} = \qty{1}{\femto\meter}
\end{equation}
Quindi:
\begin{equation}
 K_e \simeq \qty{200}{\mega\electronvolt}
\end{equation}
Possiamo dunque spiegare l'esperimento di Rutherford con le seguenti assunzioni sull'interazione:
\begin{itemize}
 \item avviene tra un elettrone e un singolo bersaglio in quiete
 \item con gli elettroni atomici è trascurabile
 \item avviene tramite la forza di Coulomb
\end{itemize}

\addfigure[Schema dell'esperimento di scattering di Rutherford]{drawio-pdf/Rutherford.drawio}{0.5}

Considerando un sottile strato di oro otteniamo il seguente risultato per la sezione d'urto differenziale:
\begin{equation}
 \dv{\sigma}{\Omega} = \brackets{\dfrac{Z_1Z_2 e^2}{4 E}}^2 \dfrac{1}{\sin^4 \brackets{\theta/2}} = \dfrac{a^4}{16}\dfrac{1}{\sin^4 \brackets{\theta/2}}
\end{equation}
Osserviamo che al variare dell'angolo di incidenza anche la sezione d'urto differenziale varia in modo significativo.\\
Possiamo inoltre identificare un parametro d'urto $b = a/2\cot(\theta/2)$, che ci permette di individuare una distanza minima $r_m$ tra la particella incidente e il nucleo.

\addtikz[Parametro d'urto e distanza minima di avvicinamento]{
    \draw[dashed,thick] (8,0) -- (0,0)
    (4,0) -- (0,2);
    \draw[-stealth,thick] (8,1) -- (5.5,1) to[out=180,in=-26.6] (2.5,1.75) -- (0,3);
    \draw[-stealth,thick] (4,0) -- (4.25,1.125)node[midway, left] {$r_m$};
    \draw[-stealth,thick] (4,0) -- (5.95,1)node[below] {$r$};
    \fill[black] (6,1) circle (0.05)node[above] {$Z_1 e$};
    \fill[white] (4,0) circle (0.3);
    \draw[thick] (4,0) circle (0.3);
    \node at (4,-0.5) {$Z_2 e$};
    \draw[-stealth,thick] (8,1) -- (7,1)node[above] {$v$};
    \draw[stealth-stealth,thick] (8,0) -- (8,1)node[midway,right] {$b$};
    \draw[thick] (3,0) arc (180:153.4:1)node[midway, left] {$\theta$};
    \draw[thick] (5,1) arc (-180:-153.4:1)node[midway, left] {$\phi$};
    \draw[thick] (5,0) arc (0:26.6:1)node[midway, right] {$\phi$};
    \node at (4,0){\large +}
}{1}{parametro_urto}

La sezione d'urto differenziale non decresce in maniera uniforme a causa della dimensione atomica. Possiamo dunque identificare massimi e minimi locali.

\addfigure[Sezione d'urto differenziale sperimentale e teorica per esperimenti di scattering. Fonte: \href{https://www.researchgate.net/publication/235601812_Elastic_and_quasi-elastic_electron_scattering_off_nuclei_with_neutron_excess}{ResearchGate}]{png/cross_section_scattering}{0.7}

Passiamo ora a ricavare la sezione d'urto di diffusione. Partiamo ricordando la regola d'oro di Fermi:
\begin{equation} \label{eq:diff_urto_1}
  W_{\text{F}} = \dfrac{2 \pi}{\hbar} \abs{M_{fi}}^2 \rho_{\text{F}}(E') = \dfrac{\dot{N}(E)}{N_a N_b} = \dfrac{\phi_a \cancel{N_b} \sigma}{N_a \cancel{N_b}} = \dfrac{\sigma v_a}{V}
\end{equation}
Dove $M_{fi}$ è l'elemento di matrice della transizione:
\begin{equation}
  M_{fi} = \bra{\psi_{\text{F}}}H_{\text{int}}\ket{\psi_i} = \int \dd{V} \psi_{\text{F}}^*(\vec{r}) H_{\text{int}}(\vec{r}) \psi_i(\vec{r})
\end{equation}
Con $H$ l'Hamiltoniana dell'interazione. La densità degli stati finali è data da:
\begin{equation} \label{eq:diff_urto_2}
  \rho_{\text{F}}(E') = \dv{n(E')}{E'}
\end{equation}
Troviamo che:
\begin{equation}
  \dd{n(k')} = \dfrac{V}{(2 \pi)^3} \dd[3]{k'} = \dfrac{V}{(2 \pi \hbar)^3}\dd{p'}
\end{equation}
Unendo la \eqref{eq:diff_urto_1} e la \eqref{eq:diff_urto_2} otteniamo:
\begin{equation}
  \dfrac{2 \pi}{\hbar} \abs{M_{fi}}^2 \dv{n(E')}{E'} = \dfrac{\sigma v_a}{V} \implies \sigma = \dfrac{2 \pi}{\hbar} \dfrac{V}{v_a} \abs{M_{fi}}^2 \dv{n(E')}{E'}
\end{equation}
In regime ultrarelativistico abbiamo $v_a \simeq c$, $E' \simeq p'c$, dunque:
\begin{equation}
  \dv{p'}{E'} \simeq \dfrac{1}{c} \implies \sigma \simeq \dfrac{2 \pi}{\hbar c} V \abs{M_{fi}}^2 \dv{n(p')}{E'} = \dfrac{2 \pi}{\hbar c} V \abs{M_{fi}}^2 \dfrac{V}{(2 \pi \hbar)^3} \dfrac{p'^2 \dd{\Omega}}{c}
\end{equation}
Ricordiamo che in questo caso l'Hamiltoniana di interazione è data dal potenziale di Coulomb:
\begin{equation}
  H_{\text{int}}(\vec{r}) = e\phi(\vec{r}) = -\dfrac{Z e^2}{r}
\end{equation}
Andando ad esprimere le funzioni d'onda iniziali e finali come onde piane otteniamo:
\begin{equation}
  \psi_i(\vec{r}) = \dfrac{1}{\sqrt{V}} \exp\brackets{i \dfrac{\vec{p}}{\hbar} \cdot \vec{r}} \quad \psi_{\text{F}}(\vec{r}) = \dfrac{1}{\sqrt{V}} \exp\brackets{i \dfrac{\vec{p}'}{\hbar} \cdot \vec{r}}
\end{equation}
Dunque l'elemento di matrice diventa:
\begin{equation}
  M_{fi} = -\dfrac{Z e^2}{V} \int \dd{\vec{r}} \dfrac{1}{r}\exp\brackets{i \dfrac{(\vec{p}-\vec{p}')}{\hbar} \cdot \vec{r}}
\end{equation}
Andando a definire il momento finale come $\vec{q} = \vec{p} - \vec{p}'$ stiamo di fatto considerando una trasformata di Fourier del potenziale di Coulomb. Ricordando che:
\begin{equation}
  \mathcal{F}(\nabla^2 \phi(\vec{r})) = -q^2 \phi(\vec{q})
\end{equation}
E $\nabla^2 \phi(\vec{r}) = -4 \pi Z e \rho(\vec{r})$ otteniamo:
\begin{equation}
  \phi(\vec{q}) = \dfrac{4 \pi }{q^2} \rho(\vec{q})
\end{equation}
Dove $\rho(\vec{q})$ è la trasformata di Fourier della densità di carica del nucleo.\\
Possiamo dunque riscrivere l'elemento di matrice come:
\begin{equation}
  M_{fi} = -\dfrac{Z e^2}{V} \dfrac{4 \pi Z e}{q^2} \rho(\vec{q}) = -\dfrac{4 \pi Z^2 e^3}{V q^2} \rho(\vec{q})
\end{equation}
Sostituendo nell'espressione della sezione d'urto otteniamo:
\begin{equation}
  \pdv{\sigma}{\Omega} = \eval{\dv{\sigma}{\Omega}}_{\text{exp}} = \eval{\dv{\sigma}{\Omega}}_{\text{Ruth}} \abs{F(q)}^2
\end{equation}
Dove $\rho(\vec{q}) = Ze F(q)$ con $F$ fattore di forma del nucleo, che tiene conto della distribuzione di carica all'interno del nucleo stesso.\\
La sezione d'urto di Rutherford che risulta:
\begin{equation}
  \eval{\dv{\sigma}{\Omega}}_{\text{Ruth}} = \brackets{\dfrac{2Z \alpha\hbar c E'}{q^2 c^2}}^2 = \brackets{\dfrac{Z \alpha\hbar c}{2 E \sin^2(\theta/2)}}^2
\end{equation}
viene dunque modificata dal fattore di forma, che dipende dalla struttura interna del nucleo.\\
Elenchiamo alcuni fattori di forma per distribuzioni di carica notevoli:
\begin{table}[H]
  \centering
  \begin{tabular}{cc}
    \hline
    \textbf{Distribuzione di carica} & \textbf{Fattore di forma} $F(\vec{q})$ \\
    \hline
    Puntiforme & $1$ \\[8pt]
    Esponenziale & $(1 + q^2/a^2\hbar^2)^{-2}$ \\[8pt]
    Gaussiana & $\exp\brackets{-q^2/2a^2 \hbar^2}$ \\[8pt]
    Sfera uniforme & $3\alpha^{-3}(\sin(\alpha) - \alpha \cos(\alpha))$ con $\alpha = qR/\hbar$ \\
    \hline
  \end{tabular}
  \caption{Fattori di forma per diverse distribuzioni di carica}
\end{table}
Sperimentalmente troviamo che la densità di carica all'interno dei nuclei è simile ad una distribuzione di Fermi:
\begin{equation}
 \rho(r) = \dfrac{\rho_0}{1 + \exp\brackets{(r-c)/a}}
\end{equation}
I parametri della distribuzione sono:
\begin{itemize}
 \item $c \simeq \qty{1.07}{\femto\meter}\,A^{1/3}$ raggio a metà altezza
 \item $a \simeq \qty{0.54}{\femto\meter}$ diffusività della superficie
\end{itemize}
Il parametro $c$ ci permette di stimare il raggio nucleare in funzione del numero di massa $A$, infatti, per una sfera uniforme, abbiamo:
\begin{equation}
    \avg{r^2} = \dfrac{\int r^2 \rho(r) \dd{V}}{\int \rho(r) \dd{V}} = \dfrac{4 \pi \int_0^{R} r^4 \rho(r) \dd{r}}{4 \pi \int_0^{R} r^2 \rho(r) \dd{r}} = \dfrac{3}{5} R^2 \implies R = \sqrt{\dfrac{5}{3}\avg{r^2}}
\end{equation}
Per la distribuzione di Fermi otteniamo:
\begin{equation}
 \sqrt{\avg{r^2}} = \sqrt{\int r^2 \rho(r) \dd{V}} \simeq r_2 A^{1/3}
\end{equation}
Con $r_2 \simeq \qty{0.94}{\femto\meter}$. Dunque:
\begin{equation}
 R \simeq \sqrt{\dfrac{5}{3}}r_2A^{1/3} = \qty{1.21}{\femto\meter} A^{1/3}
\end{equation}
Andando a valutare numericamente l'integrale osserviamo che questa approssimazione è in accordo con i dati sperimentali.

\addfigure[Raggio nucleare in funzione del numero di massa]{drawio-pdf/Nuclear_radius.drawio}{0.9}

\section{Interazioni nucleari e fissione}
Per atomi con potenziale di Coulomb trascurabile, abbiamo $N=Z$, con l'aumento del numero di massa il numero di neutroni tende ad aumentare più rapidamente rispetto al numero di protoni.\\
Considerando $A = Z+N$ il numero di massa possiamo osservare che anche la distribuzione di particelle nel nucleo assume l'aspetto di una distribuzione di Fermi:
\begin{equation}
 \rho(0) = \dfrac{A}{V} = \dfrac{3A}{4 \pi R^3} = \dfrac{3\dfrac{R^3}{r_0^3}}{4 \pi R^3} = \dfrac{3}{4 \pi r_0^3}
\end{equation}
Dunque abbiamo che $\rho \simeq \rho_0$ vicino al nucleo e decade rapidamente a zero oltre il raggio nucleare $R \simeq \qty{1}{\angstrom} $.
Possiamo anche calcolare l'energia di legame al variare del numero di massa:
\begin{equation}
 \dfrac{\text{BE}}{c^2} = \mass(A,Z) + ZM_p + N M_n
\end{equation}
Dato che $\mathcal{M}(A,Z) < ZM_p + N M_n$ l'energia di legame risulta essere negativa in quanto forza attrattiva.\\
Osserviamo che l'energia di legame per nucleone raggiunge un massimo intorno a $Fe^{56}$, per poi decrescere lentamente all'aumentare del numero di massa. Questo implica che i nuclei più leggeri tendono a fondersi, mentre quelli più pesanti tendono a fissionarsi.

\addfigure[Energia di legame per nucleone in funzione del numero di massa. Fonte: \href{https://www.britannica.com/science/nuclear-binding-energy}{Britannica}]{png/Energia_legame}{0.7}

La fissione nucleare è un processo stocastico, dunque possiamo registrare i dati sperimentali delle coppie generate. Otteniamo in questo modo una distribuzione di probabilità delle possibili coppie di frammenti.\\
Prendiamo in esame la fissione del $^{235}\atomsymbol{U}$. I prodotti hanno una $\text{BE}$ media per nucleone di circa $\qty{8.5}{\mega\electronvolt}$, mentre il $^{235}\atomsymbol{U}$ ha una $\text{BE}$ per nucleone di circa $\qty{7.6}{\mega\electronvolt}$. La differenza di energia viene rilasciata per un ammontare di circa $\qty{200}{\mega\electronvolt}$ per evento di fissione. Ad esempio:
\begin{equation}
 n + ^{235}\atomsymbol{U} \to ^{140}\atomsymbol{Xe} + ^{94}\atomsymbol{Sr} + 2n
\end{equation}
Abbiamo:
\begin{equation}
  \begin{split}
    \cancel{M_nc^2} + M_\atomsymbol{U} c^2 = M_{\atomsymbol{Xe}} c^2 + M_{\atomsymbol{Sr}} c^2 + \cancel{2}M_nc^2 + K_{\atomsymbol{Xe}} + K_{\atomsymbol{Sr}} + K_{n_1} + K_{n_2} \\[8pt]
    \underbrace{[M_\atomsymbol{U} - M_{\atomsymbol{Xe}} - M_{\atomsymbol{Sr}}- M_n]c^2}_{Q} = K_{\atomsymbol{Xe}} + K_{\atomsymbol{Sr}} + K_{n_1} + K_{n_2}
  \end{split}
\end{equation}
Andando a esprimere le masse in termini di singoli nucleoni e energie di legame otteniamo:
\begin{equation}
 Q = \text{BE}(\text{Xe}) + \text{BE}(\text{Sr}) - \text{BE}(\text{U}) + N_n \text{BE}(n)
\end{equation}
Possiamo anche calcolare il rapporto $Z/A$ del nucleo di uranio e dei frammenti di fissione:
\begin{itemize}
  \item $\dfrac{Z}{A} (\text{U}) = \dfrac{92}{235} \simeq 0.391 $
  \item $\dfrac{Z}{A}(\text{Xe}) = \dfrac{54}{140} \simeq 0.386 $
  \item $\dfrac{Z}{A}(\text{Sr}) = \dfrac{38}{94} \simeq 0.404$
\end{itemize}
I frammenti hanno dunque un rapporto $Z/A$ più vicino a quello dell'atomo iniziale, ma lontano dal valore di stabilità ottenuto da atomi con lo stesso numero di massa. Dunque i prodotti di fissione sono fortemente radioattivi e tendono a decadere per raggiungere uno stato più stabile.\\
Vogliamo ora modellizzare l'andamento della energia di legame in funzione del numero di massa. In particolare cerchiamo una combinazione di termini che possano riprodurre il grafico sperimentale:
\begin{equation}
 \text{BE}(A,Z) = \bigsum_{k=1}^{n} b_k(A,Z)
\end{equation}
Dove i termini $b_k$ rappresentano contributi diversi all'energia di legame totale. Il contributo più semplice è il termine legato al massimo valore ottenuto sperimentale. Questo dato è un termine lineare che lega l'energia di legame al numero di massa ovvero al volume del nucleo:
\begin{equation}
  R \propto A^{1/3} \implies V \propto A \implies b_1(A,Z) = -b_{\text{vol}} A
\end{equation}
Successivamente abbiamo che i nucleoni sulla superficie del nucleo hanno meno legami rispetto a quelli interni, dunque dobbiamo aggiungere un termine di correzione legato alla superficie, che deve essere positivo in quanto riduce l'energia di legame totale:
\begin{equation}
  R \propto A^{1/3} \implies S \propto A^{2/3} \implies b_2(A,Z) = b_{\text{surf}} A^{2/3}
\end{equation}
Questo termine risulta particolarmente importante per nuclei leggeri e medi $(A < 20)$.\\
Un terzo termine dipenderà dalla simmetria tra il numero di protoni e neutroni. Abbiamo che:
\begin{equation}
 b_3(A,Z) = b_{\text{sym}} \dfrac{(N-Z)^2}{A}
\end{equation}
Questo termine penalizza nuclei con un eccesso di neutroni o protoni, ovvero con numero di massa elevato, in quanto riduce l'energia di legame totale.\\
Ora dobbiamo considerare il termine di repulsione elettrostatica dovuto al potenziale coulombiano:
\begin{equation}
 b_4(A,Z) = b_{\text{coul}} \dfrac{Z^2}{A^{1/3}}
\end{equation}
Infine abbiamo un termine di accoppiamento che dipende dalla parità del numero di protoni e neutroni:
\begin{equation}
 b_5(A,Z) = \Delta P =
 \begin{cases}
  -\delta & \text{se N e Z sono pari}\\
  0 & \text{se N o Z sono dispari}\\
  +\delta & \text{se N e Z sono dispari}
 \end{cases}
\end{equation}
Sommando tutti i contributi otteniamo la formula semi-empirica della massa o di Bethe-Weizsäcker.
\begin{theorem}{Formula semi-empirica della massa (Bethe-Weizsäcker)}
    L'energia di legame di un nucleo con numero di massa $A$ e numero atomico $Z$ è data da:
    \begin{equation}
        \text{BE}(A,Z) = -b_{\text{vol}} A + b_{\text{surf}} A^{2/3} + b_{\text{sym}} \dfrac{(N-Z)^2}{A} + b_{\text{coul}} \dfrac{Z^2}{A^{1/3}} + \Delta P
    \end{equation}
    Con i seguenti valori per i coefficienti:
    \begin{itemize}
        \item $b_{\text{vol}} \simeq \qty{15.67}{\mega\electronvolt}$
        \item $b_{\text{surf}} \simeq \qty{17.23}{\mega\electronvolt}$
        \item $b_{\text{sym}} \simeq \qty{23.29}{\mega\electronvolt}$
        \item $b_{\text{coul}} \simeq \qty{0.714}{\mega\electronvolt}$
        \item $\delta \simeq 11.2 A^{-1/2}\,\text{MeV}$
    \end{itemize}
\end{theorem}
Andiamo ora a calcolare il coefficiente coulombiano. Consideriamo un nucleo come una sfera carica uniformemente distribuita con raggio $R$. La densità di carica risulta essere:
\begin{equation}
 \rho = \dfrac{Ze}{4/3 \pi R^3}
\end{equation}
Chiamiamo la carica totale $Q = Ze$. Se ipotizziamo che la densità di carica sia distribuita uniformemente all'interno del nucleo possiamo anche scrivere:
\begin{equation}
  \rho (r) = \dfrac{q}{4/3 \pi r^3} = \dfrac{Q}{4/3 \pi R^3} \implies r(q) = R \brackets{\dfrac{q}{Q}}^{1/3}
\end{equation}
Ovvero troviamo la sfera di raggio $r$ che contiene una carica $q$. Possiamo dunque calcolare l'energia potenziale elettrostatica come:
\begin{equation}
  U = \int_0^{Q} \phi(r(q)) \dd{q}
\end{equation}
Dove il potenziale all'interno della sfera carica è dato da:
\begin{equation}
 \phi(r) = \dfrac{q}{r}
\end{equation}
Sostituendo otteniamo:
\begin{equation}
  U = \int_0^{Q} \dfrac{q}{R \brackets{q/Q}^{1/3}} \dd{q} = \dfrac{Q^{1/3}}{R} \int_0^{Q} q^{2/3} \dd{q} = \dfrac{Q^{1/3}}{R}\eval{\dfrac{3}{5} q^{5/3}}_{0}^{Q} = \dfrac{3}{5} \dfrac{Q^2}{R}
\end{equation}
Risostituendo $Q = Ze$ e $R = r_0 A^{1/3}$ otteniamo:
\begin{equation}
 U = \dfrac{3}{5} \dfrac{(Ze)^2}{r_0 A^{1/3}} = b_{\text{coul}} \dfrac{Z^2}{A^{1/3}}
\end{equation}
Dunque il coefficiente coulombiano risulta essere:
\begin{equation}
 b_{\text{coul}} = \dfrac{3 e^2}{5 r_0} \simeq \qty{0.714}{\mega\electronvolt}
\end{equation}
Ora ricordiamo che la massa del nucleo è data da:
\begin{equation}
 \mathcal{M}(A,Z) = \dfrac{\text{BE}}{c^2} + ZM_p + N M_n
\end{equation}
Andiamo a trovare il valore di $Z$ che minimizza la massa del nucleo per un dato numero di massa $A$. Poniamo dunque:
\begin{equation}
 \pdv{\mathcal{M}(A,Z)}{Z} = 0
\end{equation}
Dunque:
\begin{equation}
 \begin{split}
  &\pdv{}{Z} \brackets{-b_{\text{vol}} A + b_{\text{surf}} A^{2/3} + b_{\text{sym}} \dfrac{(A-2Z)^2}{A} + b_{\text{coul}} \dfrac{Z^2}{A^{1/3}} + \Delta P + ZM_p + (A-Z) M_n} = \\[8pt]
  &-4b_{\text{sym}} \dfrac{A-2Z}{A} + 2b_{\text{coul}} \dfrac{Z}{A^{1/3}} + M_p - M_n = 0 \implies \\[8pt]
  &Z = \dfrac{A}{2} \dfrac{1 + \dfrac{M_n - M_p}{4 b_{\text{sym}}}}{1 + \dfrac{b_{\text{coul}} A^{2/3}}{4 b_{\text{sym}}}}
 \end{split}
\end{equation}
Poiché $M_n - M_p \simeq \qty{1.29}{\mega\electronvolt}$ e per valori di $A$ abbastanza grandi possiamo trascurare il termine al numeratore e approssimare $Z$ con $Z_{\text{min}}$, otteniamo:
\begin{equation}
 Z_{\text{min}} = \dfrac{A}{2} \dfrac{1}{1 + \dfrac{b_{\text{coul}} A^{2/3}}{4 b_{\text{sym}}}}
\end{equation}
Dunque se la quantità:
\begin{equation}
  \dfrac{b_{\text{coul}} A^{2/3}}{4 b_{\text{sym}}} \ll 1
\end{equation}
Abbiamo che $Z_{\text{min}} \simeq A/2$, ovvero il numero di protoni è circa uguale al numero di neutroni. Questo avviene per nuclei leggeri. Al contrario, per nuclei pesanti, il termine di repulsione elettrostatica diventa dominante e $Z_{\text{min}}$ risulta essere:
\begin{equation}
  Z_{\text{min}} \simeq \dfrac{A}{2} \dfrac{1}{\dfrac{b_{\text{coul}} A^{2/3}}{4 b_{\text{sym}}}} = \dfrac{2 b_{\text{sym}}}{b_{\text{coul}}} A^{1/3} < \dfrac{A}{2}
\end{equation}
Ovvero è significativamente minore di $A/2$. Questo spiega perché i nuclei più pesanti tendono ad avere un eccesso di neutroni rispetto ai protoni per mantenere la stabilità nucleare. Se teniamo il valore di $A$ fisso e andiamo a variare $Z$ possiamo tracciare la massa del nucleo in funzione di $Z$. Otteniamo una curva a forma di parabola con un minimo in corrispondenza di $Z_{\text{min}}$. Dunque, per un dato numero di massa $A$, il nucleo più stabile è quello con il numero di protoni che minimizza la massa totale. Ad esempio, per $A=140$ otteniamo:
\begin{equation}
  Z_{\text{min}} = \dfrac{140}{2} \dfrac{1}{1 + \dfrac{0.714 \cdot 140^{2/3}}{4 \cdot 23.29}} \simeq 57.8 \rightarrow \atom{Ce}{140}{58}
\end{equation}
Altri atomi con lo stesso numero di massa saranno instabili e tenderanno a decadere verso l'isotopo di cerio. Partendo dallo Xenon abbiamo:
\begin{equation}
  \atom{Xe}{140}{54} \to \atom{Cs}{140}{55} \to \atom{Ba}{140}{56} \to \atom{La}{140}{57} \to \atom{Ce}{140}{58}
\end{equation}
Nel tracciare la massa in funzione di $Z$ per un dato $A$ distinguiamo due casi:
\begin{itemize}
    \item $A$ pari: in questo caso abbiamo due curve distinte a seconda della parità di $N$ e $Z$ ($\pm\delta$). La curva più bassa si ha per nuclei con $N$ e $Z$ pari, mentre la curva più alta per nuclei con $N$ e $Z$ dispari. In questo caso ogni decadimento $\beta$ farà saltare da una curva all'altra fino a raggiungere il minimo di massa in corrispondenza di $Z_{\text{min}}$ \textbf{solo se} il decadimento constituisce un \textbf{vantaggio energetico}.
    \item $A$ dispari: in questo caso abbiamo una sola curva poichè uno tra $N$ e $Z$ è necessariamente dispari. Se $Z < Z_{\text{min}}$ il nucleo tenderà a decadere con decadimento $\beta^-$, mentre se $Z > Z_{\text{min}}$ il nucleo tenderà a decadere con decadimento $\beta^+$.
\end{itemize}
Ora possiamo spiegare il termine di simmetria utilizzando il modello di Fermi gas. Consideriamo un nucleo con $Z$ protoni e $N$ neutroni, entrambi trattati come gas di Fermi indipendenti. Abbiamo che:
\begin{equation}
    \begin{split}
        N = V\gamma_{\text{s}} \int \dfrac{\dd{\vec{k}}}{(2\pi)^3} f_{\text{F}}^n(k) \\[8pt]
        Z = V\gamma_{\text{s}} \int \dfrac{\dd{\vec{k}}}{(2\pi)^3} f_{\text{F}}^p(k)
    \end{split}
\end{equation}
Dove $\gamma_{\text{s}} = 2$ è il fattore di degenerazione di spin e $f_{\text{F}}^{n/p}(k)$ sono le funzioni di distribuzione di Fermi per neutroni e protoni. A temperatura nulla abbiamo:
\begin{equation}
 f_{\text{F}}^{n/p}(k) =
 \begin{cases}
  1 & k \leq k_{\text{F}}^{n/p} \\[6pt]
  0 & k > k_{\text{F}}^{n/p}
 \end{cases} \implies f_{\text{F}}^{n/p} = \Theta(k_{\text{F}}^{n/p} - k)
\end{equation}
Da cui risolviamo gli integrali ottenendo:
\begin{equation}
    \begin{split}
        N = V \dfrac{(k_{\text{F}}^n)^3}{3 \pi^2} \\[8pt]
        Z = V \dfrac{(k_{\text{F}}^p)^3}{3 \pi^2}
    \end{split}
\end{equation}
Ricordiamo che il volume del nucleo è dato da $V = 4/3 \pi R^3 = 4/3 \pi r_0^3 A$. Dunque:
\begin{equation}
    \begin{split}
        p_{\text{F}}^n = \brackets{ \dfrac{9 \pi}{4} }^{1/3} \dfrac{\hbar}{r_0} \brackets{\dfrac{N}{A}}^{1/3} \\[8pt]
        p_{\text{F}}^p = \brackets{ \dfrac{9 \pi}{4} }^{1/3} \dfrac{\hbar}{r_0} \brackets{\dfrac{Z}{A}}^{1/3}
    \end{split}
\end{equation}
Poiché nei nuclei stabili abbiamo $N \simeq Z \simeq A/2$ possiamo approssimare:
\begin{equation}
    p_{\text{F}}^n \simeq p_{\text{F}}^p \simeq p_{\text{F}} = \brackets{ \dfrac{9 \pi }{8} }^{1/3} \dfrac{\hbar}{r_0} \simeq \qty{250}{\mega\electronvolt\per c}
\end{equation}
Analogamente l'energia di Fermi risulta essere:
\begin{equation}
    E_{\text{F}} = \dfrac{(p_{\text{F}})^2}{2M} \simeq \qty{33}{\mega\electronvolt}
\end{equation}
Passiamo ora a considerare l'energia totale del sistema:
\begin{equation}
    K_{\text{tot}} = \avg{K}_n + \avg{K}_p = N \dfrac{\avg{K}_n}{N} + Z \dfrac{\avg{K}_p}{Z}
\end{equation}
Dove l'energia cinetica media per neutroni e protoni è data da:
\begin{equation}
    \begin{split}
        \dfrac{\avg{K}_{n}}{N} &= \dfrac{\displaystyle\int \dfrac{\dd{\vec{k}}}{(2\pi)^3} K_n \Theta (k_{\text{F}}^n - k)}{\displaystyle\int \dfrac{\dd{\vec{k}}}{(2\pi)^3} \Theta (k_{\text{F}}^n - k)} = \\[8pt]
        &= \dfrac{\displaystyle\int_0^{p_{\text{F}}^n} \dd{p} p^2 K_n }{\displaystyle\int_0^{p_{\text{F}}^n} \dd{p} p^2} = \dfrac{1}{2M}\dfrac{\displaystyle\int_0^{p_{\text{F}}^n} \dd{p} p^4 }{\displaystyle\int_0^{p_{\text{F}}^n} \dd{p} p^2} = \dfrac{3}{5} E_{\text{F}}^n \\[8pt]
    \end{split}
\end{equation}
Analogamente per i protoni otteniamo:
\begin{equation}
    \dfrac{\avg{K}_{p}}{Z} = \dfrac{3}{5} E_{\text{F}}^p
\end{equation}
Dunque l'energia totale del sistema risulta essere:
\begin{equation}
    K_{\text{tot}} = \dfrac{3}{5} \brackets{N E_{\text{F}}^n + Z E_{\text{F}}^p}
\end{equation}
Ricordando che $E_{\text{F}}^{n/p} \propto (N/Z)^{2/3}$ possiamo riscrivere l'energia di Fermi per neutroni e protoni come:
\begin{equation}
    \begin{split}
        E_{\text{F}}^n = 2^{2/3}E_{\text{F}} \brackets{\dfrac{N}{A}}^{2/3} \\[8pt]
        E_{\text{F}}^p = 2^{2/3} E_{\text{F}} \brackets{\dfrac{Z}{A}}^{2/3}
    \end{split}
\end{equation}
Da cui otteniamo:
\begin{equation}
    K_{\text{tot}} = \dfrac{3}{5} E_{\text{F}} 2^{2/3} \sqbr{N \brackets{\dfrac{N}{A}}^{2/3} + Z \brackets{\dfrac{Z}{A}}^{2/3}}
\end{equation}
Ora definiamo $\Delta = N-Z$, dunque:
\begin{equation}
    \begin{split}
        N &= \dfrac{A}{2}\brackets{1 + \dfrac{\Delta}{A}} \\[8pt]
        Z &= \dfrac{A}{2}\brackets{1 - \dfrac{\Delta}{A}}
    \end{split}
\end{equation}
Sviluppiamo i termini al secondo ordine in $\Delta/A$:
\begin{equation}
    \begin{split}
        N \brackets{\dfrac{N}{A}}^{2/3} &= \dfrac{A}{2^{5/3}} \brackets{1 + \dfrac{\Delta}{A}}^{5/3} \sim \dfrac{A}{2^{5/3}} \brackets{1 + \dfrac{5}{3} \dfrac{\Delta}{A} + \dfrac{5}{9} \dfrac{\Delta^2}{A^2}} \\[8pt]
        Z \brackets{\dfrac{Z}{A}}^{2/3} &= \dfrac{A}{2^{5/3}} \brackets{1 - \dfrac{\Delta}{A}}^{5/3} \sim \dfrac{A}{2^{5/3}} \brackets{1 - \dfrac{5}{3} \dfrac{\Delta}{A} + \dfrac{5}{9} \dfrac{\Delta^2}{A^2}}
    \end{split}
\end{equation}
Sostituendo otteniamo:
\begin{equation}
    K_{\text{tot}} \simeq \dfrac{3}{5} E_{\text{F}} \dfrac{A}{2^{2/3}} 2^{2/3}\sqbr{1 + \dfrac{5}{9} \dfrac{(N-Z)^2}{A^2}}
\end{equation}
Poiché l'energia di volume è:
\begin{equation}
    E_{\text{vol}} = \dfrac{3}{5} E_{\text{F}} A
\end{equation}
Otteniamo un secondo termine che dipende dalla differenza tra numero di neutroni e protoni:
\begin{equation}
    E_{\text{sym}} = \dfrac{1}{3} E_{\text{F}} \dfrac{(N-Z)^2}{A}
\end{equation}
Possiamo dunque identificare il coefficiente di simmetria come:
\begin{equation}
    b_{\text{sym}} = \dfrac{1}{3} E_{\text{F}} = \brackets{ \dfrac{9 \pi }{8} }^{2/3} \dfrac{\hbar^2}{6r_0^2 M} \simeq \qty{12}{\mega\electronvolt}
\end{equation}
Osserviamo che questo valore è circa la metà di quello ottenuto sperimentalmente, il che indica che il modello di Fermi gas, sebbene utile per comprendere qualitativamente il termine di simmetria, non cattura completamente la complessità delle interazioni nucleari reali.\\

Elenchiamo infine alcune proprietà generali delle forze nucleari:
\begin{itemize}
  \item All'interno del nucleo la forza nucleare è assai più intensa delle interazioni elettromagnetiche.
  \item La forza nucleare ha una componente attrattiva dominante.
  \item La forza nucleare ha un raggio d'azione molto limitato (circa 1-2 fm, pari alla distanza media tra i nucleoni in un nucleo).
  \item Le forze nucleari hanno proprietà di saturazione, essendo la densità centrale pressoché costante per tutti i nuclei.
  \item L'intensità della forza nucleare non può essere misurata direttamente dalle energie di legame osservate nei nuclei in quanto in esse è contenuto il temine di energia cinetica (positivo) confrontabile con il termine di energia potenziale.
  \item La forza nucleare dipende non soltanto dalla distanza relativa fra due nucleoni, ma anche da gradi di libertà intrinseci: lo spin e la carica.
  \item L'urto nucleone-nucleone a grandi energie (>200 MeV) mostra che, a distanze relative molto piccole, l'interazione diventa fortemente repulsiva.
\end{itemize}
